\chapter{Επίδειξη λειτουργίας}
\label{ch:epideixi}

Το κεφάλαιο παρουσιάζει το σύστημα σε λειτουργία, ακολουθώντας έναν χρήστη από τη σύνδεση
έως την ολοκλήρωση μιας ανταλλαγής, κατά το πρότυπο του αντίστοιχου κεφαλαίου της
\cite{hadjichristofi2022}. Όλα τα στιγμιότυπα ελήφθησαν από την εφαρμογή συνδεδεμένη σε
τοπικό δίκτυο \en{Sui}, με πραγματικές συναλλαγές. Στο σενάριο συμμετέχουν τρεις
διευθύνσεις: ο χρήστης που παρακολουθούμε, του οποίου η διεύθυνση εμφανίζεται επάνω δεξιά
σε κάθε στιγμιότυπο μετά τη σύνδεση, και δύο αντισυμβαλλόμενοι που κατέχουν \en{NFTs} και υποβάλλουν
προσφορές. Τα στιγμιότυπα παρήχθησαν από το σενάριο
αυτοματοποιημένης πλοήγησης της Ενότητας~\ref{sec:verification}, ώστε να αναπαράγονται
πανομοιότυπα μετά από κάθε αλλαγή της διεπαφής.

\newcommand{\shot}[2][\textwidth]{{\setlength{\fboxsep}{0pt}\setlength{\fboxrule}{0.5pt}%
  \fcolorbox{gray!55}{white}{\includegraphics[width=\dimexpr#1-1pt\relax]{shots/#2}}}}
\newcommand{\shotpair}[2]{\shot[0.49\textwidth]{#1}\hfill\shot[0.49\textwidth]{#2}}
\newcommand{\shottriple}[3]{\shot[0.31\textwidth]{#1}\hfill\shot[0.31\textwidth]{#2}\hfill\shot[0.31\textwidth]{#3}}

\section{Αρχική σελίδα και σύνδεση πορτοφολιού}
\label{sec:demo-connect}

Η αρχική σελίδα (Σχήμα~\ref{fig:demo-home}) εξηγεί σε μία πρόταση τι κάνει η εφαρμογή,
συνοψίζει τη διαδικασία σε τρία βήματα και παρουσιάζει τις τρεις συλλογές με το πλήθος των
\en{NFTs} που έχουν δημιουργηθεί σε καθεμία. Είναι προσβάσιμη χωρίς πορτοφόλι.

\begin{figure}[H]
\centering
\shot{01-home-disconnected.jpg}
\caption{Η αρχική σελίδα πριν από τη σύνδεση πορτοφολιού.}
\label{fig:demo-home}
\end{figure}

Με το κουμπί «\en{Connect Wallet}» ανοίγει το παράθυρο επιλογής πορτοφολιού
(Σχήμα~\ref{fig:demo-connect}, αριστερά), όπου εμφανίζονται όσα πορτοφόλια υλοποιούν το
\en{Wallet Standard}. Μετά τη σύνδεση, το κουμπί δείχνει τη διεύθυνση του χρήστη (δεξιά),
δίπλα του εμφανίζονται το υπόλοιπο \en{SUI} και το κουμπί «\en{Faucet}» για δοκιμαστικά
νομίσματα, και οι σελίδες «\en{Mint}», «\en{My NFTs}» και «\en{My bids}» ενεργοποιούνται.

\begin{figure}[H]
\centering
\shotpair{02-connect-modal.jpg}{08-home-connected.jpg}
\caption{Αριστερά: επιλογή πορτοφολιού. Δεξιά: η αρχική σελίδα με συνδεδεμένο πορτοφόλι
και \en{NFTs} ήδη στην αγορά.}
\label{fig:demo-connect}
\end{figure}

\section{Δημιουργία NFT}
\label{sec:demo-mint}

Η σελίδα «\en{Mint}» (Σχήμα~\ref{fig:demo-mint}) παρουσιάζει τα εννέα είδη του καταλόγου,
τρία ανά συλλογή, με τη βαθμίδα σπανιότητάς τους. Ο χρήστης ορίζει ποσότητα και πατά
«\en{Mint}».

\begin{figure}[H]
\centering
\shot{03-mint.jpg}
\caption{Η σελίδα δημιουργίας \en{NFT}.}
\label{fig:demo-mint}
\end{figure}

Όσο η συναλλαγή εκκρεμεί (Σχήμα~\ref{fig:demo-mint-states}, αριστερά), η κάρτα που την
εκκίνησε επισημαίνεται και το κουμπί της δείχνει «\en{Minting…}». Οι υπόλοιπες κάρτες
παραμένουν ενεργές: ένα πάτημα σε άλλη κάρτα μπαίνει σε ουρά και εκτελείται μόλις
ολοκληρωθεί η τρέχουσα συναλλαγή. Με την ολοκλήρωση εμφανίζεται ειδοποίηση (δεξιά) με
σύνδεσμο προς τη συναλλαγή στον εξερευνητή, και το υπόλοιπο στη γραμμή πλοήγησης
ενημερώνεται. Στο παράδειγμα δημιουργήθηκε ένα \en{Common Sedan}.

\begin{figure}[H]
\centering
\shotpair{04-mint-pending.jpg}{05-mint-toast.jpg}
\caption{Αριστερά: εκκρεμής δημιουργία. Δεξιά: ειδοποίηση επιτυχίας.}
\label{fig:demo-mint-states}
\end{figure}

\section{Περιήγηση σε συλλογή}
\label{sec:demo-collection}

Επιλέγοντας μια συλλογή από την αρχική σελίδα, ο χρήστης βλέπει όλα τα \en{NFTs} της που
βρίσκονται σε πορτοφόλια (Σχήμα~\ref{fig:demo-collection}). Κάθε κάρτα δείχνει την εικόνα,
τη σπανιότητα ως μεταλλική ετικέτα, το όνομα, και τη διεύθυνση του κατόχου με σύνδεσμο προς
\en{explorer}. Τα \en{NFTs} άλλων χρηστών φέρουν το κουμπί «\en{Place a bid}». Το \en{NFT}
που ανήκει στον ίδιο τον χρήστη διακρίνεται με περίγραμμα και την ένδειξη
«\en{You own this}». Τα τρία πλήκτρα κάτω από τον τίτλο φιλτράρουν τη συλλογή ανά
βαθμίδα σπανιότητας, με το πλήθος κάθε βαθμίδας, και η επιλογή αποτυπώνεται στη διεύθυνση
της σελίδας ώστε να μπορεί να μοιρασθεί.

\begin{figure}[H]
\centering
\shot{09-collection.jpg}
\caption{Η σελίδα της συλλογής \en{Cars}. Πρώτη αριστερά η κάρτα \en{NFT} του ίδιου του
χρήστη.}
\label{fig:demo-collection}
\end{figure}

\section{Υποβολή προσφοράς}
\label{sec:demo-bid}

Το «\en{Place a bid}» ανοίγει τη σύνθεση προσφοράς σε πλήρη οθόνη. Η αριστερή στήλη δείχνει
σταθερά τι θα \emph{λάβει} ο χρήστης· η δεξιά, τι \emph{δίνει}. Η γραμμή προόδου στο επάνω
μέρος δείχνει τα τρία βήματα.

Στο πρώτο βήμα (Σχήμα~\ref{fig:demo-bid-13}, αριστερά) ο χρήστης ορίζει προαιρετικά ένα
ποσό \en{SUI}, είτε πληκτρολογώντας το είτε με ένα από τα προτεινόμενα ποσά, ή παραλείπει
το βήμα με το «\en{Skip}». Στο δεύτερο (Σχήμα~\ref{fig:demo-bid-2}) επιλέγει \en{NFTs} από το πορτοφόλι
του. Στο παράδειγμα προσφέρονται 1{,}5 \en{SUI}, ένα \en{Uncommon Townhouse} και ένα
\en{Common Sedan} για ένα \en{Rare Supercar}: μια προσφορά που συνδυάζει νόμισμα με
\en{NFTs} δύο διαφορετικών συλλογών για ένα \en{NFT} τρίτης βαθμίδας. Η σύνοψη στο κάτω
μέρος της οθόνης ενημερώνεται με κάθε επιλογή.

\begin{figure}[H]
\centering
\shot{11-bid-step2.jpg}
\caption{Σύνθεση προσφοράς, βήμα 2: επιλογή \en{NFTs}. Κάτω αριστερά η σύνοψη
«δίνεις \swap{} παίρνεις».}
\label{fig:demo-bid-2}
\end{figure}

Στο τρίτο βήμα (Σχήμα~\ref{fig:demo-bid-13}, δεξιά) ο χρήστης βλέπει την τελική μορφή της
ανταλλαγής και μια υπενθύμιση ότι τα στοιχεία του θα δεσμευθούν έως ότου αποφασίσει ο
κάτοχος. Με το «\en{Place bid}» το πορτοφόλι ζητά έγκριση.

\begin{figure}[H]
\centering
\shotpair{10-bid-step1.jpg}{12-bid-step3.jpg}
\caption{Αριστερά: βήμα 1, ποσό \en{SUI}. Δεξιά: βήμα 3, επισκόπηση.}
\label{fig:demo-bid-13}
\end{figure}

Μετά την επιτυχή υποβολή ο χρήστης επιστρέφει στη συλλογή με σχετική ειδοποίηση, η
οποία, όπως κάθε ειδοποίηση επιτυχούς συναλλαγής, περιέχει σύνδεσμο προς τη συναλλαγή στον
εξερευνητή του τρέχοντος δικτύου. Η
προσφορά εμφανίζεται στην καρτέλα «\en{Placed}» της σελίδας «\en{My bids}»
(Σχήμα~\ref{fig:demo-placed}). Η γραμμή συνοψίζει την ανταλλαγή, δύο \en{NFTs} και 1{,}5 \en{SUI} για ένα \en{NFT}, και,
όταν αναπτυχθεί, την αναλύει στα δύο πλαίσια. Από εδώ ο
χρήστης μπορεί να την ακυρώσει.

\begin{figure}[H]
\centering
\shotpair{13-bid-placed-toast.jpg}{14-my-bids-placed.jpg}
\caption{Αριστερά: ειδοποίηση επιτυχούς υποβολής, πάνω από τη συλλογή. Δεξιά: η προσφορά
στην καρτέλα «\en{Placed}».}
\label{fig:demo-placed}
\end{figure}

\section{Εισερχόμενες προσφορές}
\label{sec:demo-incoming}

Η σελίδα «\en{My NFTs}» (Σχήμα~\ref{fig:demo-mynfts}) δείχνει τα \en{NFTs} του χρήστη. Όσα
έχουν λάβει προσφορές φέρουν ένδειξη με το πλήθος τους, ώστε ο χρήστης να βλέπει με μια
ματιά πού υπάρχει ενδιαφέρον. Όσα έχει δεσμεύσει ο ίδιος σε προσφορά του εμφανίζονται
επίσης, αχνότερα και με την ένδειξη «\en{In a bid}»· πατώντας τα βλέπει την προσφορά που
τα κρατά και μπορεί να την ακυρώσει από εκεί. Η σπανιότητα αποδίδεται εδώ με τη μορφή
κουκκίδας, που διαβάζεται και σε μικρό μέγεθος.

\begin{figure}[H]
\centering
\shot{06-my-nfts.jpg}
\caption{Τα \en{NFTs} του χρήστη: τα δύο πρώτα με ένδειξη εισερχόμενων προσφορών, τα δύο
επόμενα δεσμευμένα στη δική του προσφορά («\en{In a bid}»).}
\label{fig:demo-mynfts}
\end{figure}

Πατώντας σε ένα \en{NFT} ανοίγει το παράθυρο με τις προσφορές του
(Σχήμα~\ref{fig:demo-incoming}). Κάθε προσφορά παρουσιάζεται από την οπτική του κατόχου:
τι θα λάβει, με μικρογραφίες των προσφερόμενων \en{NFTs}, και ότι θα δώσει το συγκεκριμένο
\en{NFT}. Στο παράδειγμα το \en{Rare Mansion} έχει δύο προσφορές από διαφορετικούς
προσφέροντες: μία με ένα \en{NFT} μόνο (ένα \en{Rare Supercar}) και μία με ένα \en{NFT}
και 0{,}5 \en{SUI} (ένα \en{Uncommon Coupe}).

\begin{figure}[H]
\centering
\shot{07-incoming-bids.jpg}
\caption{Οι εισερχόμενες προσφορές για το \en{Rare Mansion}.}
\label{fig:demo-incoming}
\end{figure}

\section{Αποδοχή, απόρριψη και ακύρωση}
\label{sec:demo-settle}

Όλες οι εισερχόμενες προσφορές του χρήστη, ανεξαρτήτως \en{NFT}, συγκεντρώνονται στην
καρτέλα «\en{Received}» της σελίδας «\en{My bids}», εδώ τρεις: δύο για το \en{Rare
Mansion} και μία για το \en{Common Cottage}. Αναπτύσσοντας μια γραμμή
(Σχήμα~\ref{fig:demo-received}) ο χρήστης βλέπει την ανταλλαγή αναλυτικά: αριστερά το
\en{NFT} που θα δώσει, δεξιά το \en{SUI} και τα \en{NFTs} που θα λάβει, το καθένα με τη
σπανιότητα και το \en{ID} του.

\begin{figure}[H]
\centering
\shot{16-my-bids-received-expanded.jpg}
\caption{Η καρτέλα «\en{Received}» με τρεις εισερχόμενες προσφορές, η τρίτη ανεπτυγμένη.}
\label{fig:demo-received}
\end{figure}

Με το «\en{Accept trade}» και την έγκριση στο πορτοφόλι, η ανταλλαγή εκτελείται. Όσο
εκκρεμεί (Σχήμα~\ref{fig:demo-accept}, αριστερά), το κουμπί δείχνει «\en{Accepting…}» και
οι υπόλοιπες προσφορές απενεργοποιούνται, ώστε να μην υποβληθούν ανταγωνιστικές ενέργειες.
Με την ολοκλήρωση (δεξιά), η ειδοποίηση αναφέρει τι προστέθηκε στο πορτοφόλι, το
υπόλοιπο αυξάνεται κατά τα 0{,}5 \en{SUI} της προσφοράς, και η λίστα ανανεώνεται. Αξίζει να
παρατηρηθεί ότι από τη λίστα αφαιρέθηκε όχι μόνο η προσφορά που έγινε αποδεκτή αλλά και η
δεύτερη προσφορά για το ίδιο \en{NFT}: το \en{NFT} άλλαξε κάτοχο, και η προσφορά εκείνη
απευθύνεται πλέον στον νέο. Απομένει μόνο η προσφορά για το \en{Common Cottage}.

\begin{figure}[H]
\centering
\shotpair{17-accepting.jpg}{18-trade-complete.jpg}
\caption{Αριστερά: αποδοχή σε εξέλιξη. Δεξιά: η ανταλλαγή ολοκληρώθηκε.}
\label{fig:demo-accept}
\end{figure}

Η απόρριψη μιας προσφοράς και η ακύρωση μιας προσφοράς που έχει υποβάλει ο ίδιος ο χρήστης
ακολουθούν το ίδιο πρότυπο: ένα πάτημα, μία έγκριση, και ειδοποίηση που δηλώνει πού
κατέληξαν τα στοιχεία (Σχήμα~\ref{fig:demo-reject-cancel}). Στο παράδειγμα απορρίπτεται η
προσφορά για το \en{Common Cottage}, οπότε η καρτέλα «\en{Received}» αδειάζει, και έπειτα
ακυρώνεται η προσφορά του ίδιου του χρήστη για το \en{Rare Supercar}: τα δύο \en{NFTs}
επιστρέφουν στο πορτοφόλι του και το υπόλοιπο ανεβαίνει κατά το 1{,}5 \en{SUI} που είχε
δεσμεύσει.

\begin{figure}[H]
\centering
\shotpair{19-bid-rejected.jpg}{20-bid-cancelled.jpg}
\caption{Αριστερά: απόρριψη εισερχόμενης προσφοράς. Δεξιά: ακύρωση υποβληθείσας
προσφοράς.}
\label{fig:demo-reject-cancel}
\end{figure}

\section{Η εφαρμογή σε κινητή οθόνη}
\label{sec:demo-mobile}

Η διεπαφή προσαρμόζεται σε πλάτος 360 εικονοστοιχείων χωρίς απώλεια λειτουργιών· τα
στιγμιότυπα που ακολουθούν ελήφθησαν σε πλάτος 390. Η πλοήγηση μεταφέρεται σε δεύτερη
γραμμή και το κουμπί του \en{faucet} γίνεται εικονίδιο. Το Σχήμα~\ref{fig:demo-mobile-1}
δείχνει την αρχική σελίδα, μια συλλογή σε δίστηλο πλέγμα με τα φίλτρα σπανιότητας να
κυλούν οριζόντια, και την επισκόπηση προσφοράς. Στη στενή οθόνη τα δύο πλαίσια της
ανταλλαγής στοιβάζονται κατακόρυφα, και το κουμπί υποβολής καταλαμβάνει όλο το πλάτος.

\begin{figure}[H]
\centering
\shottriple{m1-home.jpg}{m2-collection.jpg}{m5-bid-review.jpg}
\caption{Κινητή οθόνη: αρχική σελίδα, συλλογή, επισκόπηση προσφοράς.}
\label{fig:demo-mobile-1}
\end{figure}

Το Σχήμα~\ref{fig:demo-mobile-2} δείχνει τη σελίδα δημιουργίας, όπου οι συλλογές
επιλέγονται με ετικέτες αντί να εμφανίζονται όλες μαζί· τις εισερχόμενες προσφορές, που
παρουσιάζονται ως φύλλο που ανασύρεται από το κάτω μέρος της οθόνης· και τη σελίδα
προσφορών, όπου κάθε γραμμή συνοψίζει την ανταλλαγή και τα κουμπιά ενεργειών εμφανίζονται
μόνο όταν αναπτυχθεί, σε μέγεθος κατάλληλο για αφή.

\begin{figure}[H]
\centering
\shottriple{m6-mint.jpg}{m8-incoming-bids.jpg}{m9-my-bids.jpg}
\caption{Κινητή οθόνη: δημιουργία, εισερχόμενες προσφορές, σελίδα προσφορών.}
\label{fig:demo-mobile-2}
\end{figure}
